\documentclass[12pt]{article}

\usepackage[english]{babel}
\usepackage[a4paper,top=1.5cm,bottom=1.5cm,left=2cm,right=2cm,marginparwidth=1.75cm]{geometry}
\usepackage{graphicx}
\usepackage{amsfonts}
\usepackage{amsmath}
\usepackage{amssymb}
\usepackage{enumitem}
\usepackage{booktabs}

\title{The Dominant Pole Approximation}
\author{Allan Wu}
\date{\today}

\begin{document}

\maketitle{}

\section*{Introduction and Revision}

We give a brief overview of the following topics:
\begin{enumerate}
\item The standard second-order transfer function
\item Underdamped systems and examples
\item Higher-order systems and examples
\item The dominant pole approximation (concept)
\end{enumerate}

\section*{Demonstration and Analysis}

\subsection*{Transfer Functions}

We consider a 5th order underdamped system characterised by the transfer function:
\begin{align*}
T(s) &= \frac{12614(s+5)(s+6-j8)(s+6+j8)}{5(s+7)(s+5-j9)(s+5+j9)(s+1-j4)(s+1+j4)} \\
&= \frac{12614 s^3 + 214438 s^2 + 2018240 s + 6307000}{5 s^5 + 95 s^4 + 1135 s^3 + 6915 s^2 + 22380 s + 63070}
\end{align*}

Using partial fraction expansion:
\begin{equation}
T(s) = \mathcal{L}^{-1} \left\{ \frac{10.6}{s+7}+\frac{3.747 s-38.1656}{s^{2}+10 s+106}-\frac{114.3472 s+78.8565}{s^{2}+2 s+17}+\frac{100}{s} \right\}
\end{equation}

After applying the inverse Laplace transform, the expression for the transient response is given by:
\begin{align*}
\left[ 10.6e^{-7t}+e^{-5t}\left(3.747\cos9t-6.322\sin9t\right)-e^{-t}\left(114.3472\cos4t-8.873\sin4t\right)+100 \right] u(t)
\end{align*}

Observing that all poles have real parts at least 5 times greater than the pair of dominant poles closest to the $j\omega$-axis, we can apply the simplification procedure.
\begin{enumerate}
\item Include the pair of dominant poles closest to the imaginary axis (i.e. the smallest real part) in the denominator polynomial. Remove the other poles.
\item Match the DC gain of the original transfer function. That is, multiply the numerator by an appropriate constant so that $T\left(0\right)=T_{dp}\left(0\right).$
\end{enumerate}

By applying the dominant pole approximation, we obtain a simplified 2nd order transfer function:
\begin{align*}
T_{dp}(s) &= \frac{1700}{(s+1-j4)(s+1+j4)} \\
&= \frac{1700}{s^2 + 2s + 17}
\end{align*}

\subsection*{Time Domain Analysis}

\begin{figure}
\centering
\includegraphics[width=0.9\linewidth]{stepplot.png}
\caption{Annotated step response of the original and simplified transfer functions.}
\label{fig:stepplot}
\end{figure}

\subsection*{S-Domain Analysis}
\begin{figure}
\centering
\includegraphics[width=0.6\linewidth]{pzplot.png}
\caption{Poles and zeros of the transfer function.}
\label{fig:pzplot}
\end{figure}

\subsection*{Comparative Analysis and Discussion}

TABLE 3. S-DOMAIN PARAMETERS
\begin{table}[h]
\centering
\caption{S-DOMAIN PARAMETERS.}
\label{tab:sparams}
\begin{tabular}{cccc}\toprule
Pole        &Natural frequency ${\omega }_{n}$ &Damped frequency ${\omega }_{d}$ &Damping ratio $\zeta$ \\ \midrule
$-1\pm j4$  &$2\pi\cdot 0.656$	&$2\pi\cdot 0.637$   &$0.243$ \\
$-5\pm j9$  &$2\pi\cdot 1.639$	&$2\pi\cdot 1.432$   &$0.486$ \\
$-7+j0$	    &$2\pi\cdot 1.114$	&$0 \left(DC\right)$	&$1.000$ \\
\bottomrule
\end{tabular}
\end{table}

\section*{References}
\begin{enumerate}[label={[\arabic*]}]
\item Dorf, R.C. and Bishop, R.H. (2005) \textit{Modern Control Systems} (13th ed). Section 5.4 ``Effects of a Third Pole and a Zero on the Second-Order System Response''
\item Ogata, K. (2010) \textit{Modern Control Engineering} (5th ed). Section 5.4 "Higher-Order Systems"
\end{enumerate}

\end{document}